\documentclass{article}
\usepackage{amsmath}
\usepackage{amsthm}

\newtheorem{theorem}{Problem}

\begin{document}

\begin{theorem}
How many positive integers leave a remainder of 2 when divided by 8 and are less than 100?
\end{theorem}

\begin{proof}
Step 0: If a number leaves a remainder of 2 when divided by 8, it is of the form 8n + 2.
Step 1: We need to determine how many numbers of this form are positive and less than 100.
Step 2: Start by testing small values of n.
Step 3: For n = 1, 8*1 + 2 = 10.
Step 4: For n = 2, 8*2 + 2 = 18, and so forth.
Step 5: Continue this process to find all numbers until 8n + 2 becomes 100 or more.
Step 6: For n = 12, 8*12 + 2 = 98, and for n = 13, 8*13 + 2 = 106 which is not under 100.
Step 7: The range of n that satisfies this is 1 to 12.
\end{proof}

\end{document}
